\section{Introduction}\label{sec:introduction}

National grid operators must anticipate regional load and assess system pressure on a day-ahead horizon to coordinate generation dispatch, reserve planning, and constraint management. In developing economies, supply limitations, infrastructure stress, and load shedding often coincide with rapid demand growth, so planners need forecasts that show both where load will concentrate across regions and how severely operational constraints may bind system-wide. Many South Asian systems report daily division-level balances, national generation scalars, and limitation indicators that nevertheless encode stress when reserves tighten or infrastructure constraints bind. On the Bangladesh national power system, operational data are reported across nine administrative divisions over a multi-year timeline, with regional evening-peak demand and limitation dynamics co-visible in each daily record \cite{Hossain2025BangladeshSmartGrid}. Effective planning therefore requires spatially disaggregated, temporally aligned forecasts that support review when stress indicators and load trajectories must be judged together.

Despite these data, forecasting practice often treats load prediction and stress monitoring separately. Classical and ensemble models \cite{Breiman2001RandomForest,Chen2016XGBoost,Musbah2025LoadForecasting} forecast demand without explicit inter-division graph structure, while shallow recurrent or graph-convolution architectures \cite{Hochreiter1997LSTM,Yu2018STGCN} may fail to combine long-range temporal context with informative multi-node representations on national footprints. Graph-based methods \cite{Kipf2017GCN,Wu2023TransferLearningGNN,Shi2024EVChargingGNN} from residential or microgrid settings often rely on geographical adjacency that may not capture data-driven co-movement among regional loads. Operational stress---from reserve shortfall, limitation pressure, and shedding intensity---is seldom forecast jointly with regional demand; shedding studies \cite{Paphitis2025SpatialLoadShedding,Wang2026SHAPLoadShedding} typically address control and optimisation rather than day-ahead graded stress indices. Explainability, where present \cite{Baur2024XAIReview,Cifci2025InterpretableSmartGrid}, is rarely integrated with multi-node graph-temporal models that emit coupled load and stress outputs. Operational records thus offer synchronised regional load and constraint context, whereas many pipelines deliver disconnected single-task predictions with limited spatial structure and transparency.

We address this mismatch by formulating day-ahead forecasting as a supervised multi-task problem \cite{Caruana1997Multitask} over nine spatial units. The primary task is next-day regional evening-peak demand; the secondary task is a graph-level Operational Stress Index (OSI) summarising system-wide pressure from shedding intensity, generation reserve margin, and operational limitation components normalised under a leakage-safe training policy. OSI is a bounded composite score rather than a sparse per-division shedding label, preserving graded stress magnitude when reserve or limitation pressure is elevated even without shedding. Demand and OSI are operationally related yet statistically non-redundant: demand describes expected load by region, whereas OSI encodes constraint-related pressure not fully determined by demand alone. The research problem is whether a single spatio-temporal framework can learn coupled regional demand and operational stress forecasts from historical Bangladesh grid records under chronological evaluation, and which design choices---graph prior, architecture, task weighting, and post-hoc attribution---are empirically justified on this case.

A structured literature review identifies further gaps motivating this work. Bangladesh division-level smart-grid daily data remain rare in peer-reviewed spatio-temporal load and shedding studies; developing-economy national grids are under-represented relative to US and European benchmarks. Multi-node demand forecasting is seldom coupled with operational-stress targets in one multi-task framework, and shedding research rarely targets forecast-centric graded daily stress. Only a minority of recent work combines graph-based spatial coupling with transformer-style temporal modelling \cite{Zhao2024GraphAttentionTransformer,Bentsen2023GraphTransformer,Vaswani2017Attention} for multi-node national forecasting, and none address the Bangladesh shedding and limitation context directly. Explainable artificial intelligence appears in few graph-temporal energy studies \cite{Baur2024XAIReview,Lundberg2017SHAP} and is seldom reported with multi-method discordance for coupled demand and stress outputs. Weak documentation of chronological splits, baselines, and ablations in part of the conference literature further complicates fair comparison. These gaps motivate a reproducible, evidence-locked study that couples graph-temporal learning, multi-task stress forecasting, and operator-oriented attribution within one frozen evaluation programme.

This paper proposes the Correlation-Aware Multi-Task Forecasting Framework, implemented as a Parallel-Fusion Spatio-Temporal Graph Transformer (PF-STGT). The framework ingests a seven-day lookback of node-level and national context features with a train-estimated correlation graph encoding inter-regional demand coupling from training data only. A parallel dual-branch encoder combines graph-aware spatial attention with temporal transformer capacity \cite{Vaswani2017Attention,Kipf2017GCN}; learned fusion feeds task-specific heads that decode regional megawatt demand and a bounded system-wide OSI scalar for the next operating day. Graph construction compares geographical, hybrid, and correlation-only priors; the final configuration adopts correlation-thresholded adjacency following empirical comparison. Multi-task training uses scale-aware joint loss with explicit stress weighting \cite{Caruana1997Multitask,Song2024MultiTaskEnergy} so shared representations retain limitation-relevant structure without collapsing to demand-only learning. Post-hoc attribution applies grouped integrated gradients, permutation importance, and spatial and temporal attention export \cite{Sundararajan2017IntegratedGradients,Breiman2001RandomForest,Lundberg2017SHAP}, together with stratified case studies on the frozen checkpoint, with method disagreement reported explicitly. Evaluation on chronologically partitioned data includes classical, recurrent, and graph-convolution benchmarks \cite{Breiman2001RandomForest,Hochreiter1997LSTM,Zhao2020TGCN}, configuration ablations, architecture simplification, error geography, and explainability analysis; protocol details appear in later sections.

The principal contributions are as follows. First, a Correlation-Aware Multi-Task Forecasting Framework (configuration S2) for joint day-ahead regional demand and graph-level OSI prediction from a single spatio-temporal graph-transformer on Bangladesh national grid data. Second, evidence-driven architecture selection that replaces the initial hybrid geographical--correlation graph with correlation-only adjacency while retaining full temporal encoder capacity where ablation shows it is consequential. Third, a repaired multi-task training protocol with explicit demand--stress loss balancing under the frozen W20 specification. Fourth, an empirical evaluation programme---external benchmarks, component ablations, inferential testing on primary demand MAE with multiple-comparison correction, and metric verification for classical baselines---supporting reproducible comparison within the frozen suite. Fifth, integrated explainability analysis with multi-method attribution artefacts for both tasks, including spatial and temporal views and comparison against OSI component decomposition on audited case studies. These contributions rest on offline evaluation on held-out test data and are positioned as planning and transparency evidence rather than proof of deployment outcomes or universal architectural superiority.

The remainder of this paper is organised as follows. Section~\ref{sec:related_work} reviews related work; Section~\ref{sec:methodology} presents the methodology; Section~\ref{sec:experimental_setup} describes the experimental setup; Section~\ref{sec:results} reports results; Section~\ref{sec:discussion} discusses findings and limitations; and Section~\ref{sec:conclusion} concludes. Supplementary materials are provided in Appendix~A.
